\clearpage
\part*{Part XVI. The Scorecard, The Open Results And What Remains}
\label{atlaspart:16-scorecard}
\addcontentsline{toc}{part}{Part XVI. The Scorecard, The Open Results And What Remains}
\section*{Introduction to Part XVI}\label{narr:part:XIV}
\addcontentsline{toc}{section}{Introduction to Part XVI}
\par\noindent\fbox{\parbox{0.96\linewidth}{\small Table~\ref{box:XVI-1-the-results-table-re-tallied-at-the-cut:1} counts the boxed statements of this release, \atlasTheorems{} theorems, \atlasConditionals{} conditional theorems and \atlasOpen{} open results, each listed in \texttt{compendium/index}. The other tallies of this Part also count superseded and re-posed results, so they are larger than the counts of this release.}}\par\medskip
The last Part is the ledger of the compendium. It re-tallies the results, counts the open results by form and by Part, prints the page of what the construction has not reached, and concludes. Every chapter of the Part is a table or a census, and the reader should use it as one uses an index: to find, for any claim of the compendium, where it stands and what still parts.

The status matrix beneath this introduction is the register matrix of the framework chapter read for standing. At the cell (3,2) the photon stands as the first conjecture of the compendium, attached to the open link whose office is zero and whose center acts. At the cell (2,2) the matrix reads spacetime, conjectured, and at the four remaining cells the sector names are correspondences by name, each with the count of results that support the identification and the count that part from it. The matrix is the shortest summary of how far the construction has reached across the six genres of the link, and the chapters of this Part spell out the same standing result by result.

Chapter 14-1 re-tallies the results of this release by standing. Chapter 14-2 prints the census of the open results by form, status and Part, Table~238.1, which is the table to consult for every box of this compendium that is an open result.

Chapter 14-7 is the ungraded page. Table~239.1 prints the register of what the construction has not reached, and Open result~207.1, the strong reading as a theorem of the rate law, stands open. The conclusion is held in Table~240.1.

In the language of the literature, this Part names the Standard Model and the open problems of physics. The correspondence is by name at each. The boxes are tables, which record standing and not measurement, and the standing of every result is the standing its own chapter gives it.
\label{narr:open:14-1}The ledger opens with the results of this release by standing.
\label{narr:open:14-2}Table~238.1 prints the census of the open results by form, status and Part, the table to consult for every box of this compendium that is an open result.
\label{narr:open:14-7}Table~239.1 prints the register of what the construction has not reached, and Open result~207.1, the strong reading as a theorem of the rate law, stands open.
\label{narr:open:14-C}Table~240.1 holds the conclusion, resting on the results table and the census of the open results.
\begin{figure}[htbp]\centering
\footnotesize
\setlength{\tabcolsep}{3pt}
\begin{adjustbox}{max width=\linewidth}\begin{tabular}{|p{14mm}|p{44mm}|p{44mm}|p{44mm}|}
\hline
 & \textbf{register two} (the qubit) & \textbf{register three} (the qutrit) & \textbf{register four} (the ququart) \\ \hline
\textbf{register two} & \textbf{(2,2)} \newline a conjecture \newline the pair ket at the join has a zero office and no acting center \newline \emph{spacetime, conjectured} & \textbf{(3,2)} \newline a conjecture named at a theorem \newline the open link meets both criteria, its office zero and its center acting \newline \emph{electromagnetism: the photon (Conjecture 1)} & \textbf{(4,2)} \newline a correspondence by name \newline 38 supporting results, 14 parting \newline \emph{the weak interaction, by name} \\ \hline
\textbf{register three} & \emph{the link (3,2), read from its other end} & \textbf{(3,3)} \newline a correspondence by name \newline 45 supporting results, 32 parting \newline \emph{the strong interaction, by name} & \textbf{(4,3)} \newline a correspondence by name \newline 16 supporting results, 30 parting \newline \emph{the Yukawa sector, by name} \\ \hline
\textbf{register four} & \emph{the link (4,2), read from its other end} & \emph{the link (4,3), read from its other end} & \textbf{(4,4)} \newline a correspondence by name \newline 25 supporting results, 20 parting \newline \emph{the Higgs, by name} \\ \hline
\end{tabular}\end{adjustbox}
\caption{The status matrix: the standing of every cell of the register matrix at this release of the compendium. At the four cells named by correspondence the correspondence is by name. At the cell (3,2) the photon is the first conjecture of the compendium, attached to the open link whose office is zero and whose center acts, and at the cell (2,2) the matrix reads spacetime, conjectured.}\label{fig:register-matrix-status}
\end{figure}
The face group of every cell, read from the group row of the compendium, is the special unitary pair group of its two registers, and the six names the literature gives these groups are, cell by cell, at (2,2) Spin$^{c}$(4), at (2,3) the Standard Model group, at (2,4) S(U(2) x U(4)) with SU(4) $=$ Spin(6), at (3,3) S(U(3) x U(3)), at (3,4) S(U(3) x U(4)), at (4,4) S(U(4) x U(4)); the correspondence is by name at every cell where the status matrix says so.
\section{The results table re-tallied}\label{mod:XVI-1-the-results-table-re-tallied-at-the-cut}
\chdesig{16-1}{}{XVI-1-the-results-table-re-tallied-at-the-cut}{XVI-1-the-results-table-re-tallied-at-the-cut}
\label{narr:open:14-1}The ledger opens with the results of this release by standing.
\begin{chtable}[The results of this compendium by standing]\label{box:XVI-1-the-results-table-re-tallied-at-the-cut:1}
This release holds \atlasTheorems{} theorems, \atlasConditionals{} conditional theorems and \atlasOpen{} open results among its boxed statements, and \atlasTheoremsLean{} of the \atlasTheorems{} theorems are tied to named Lean declarations, which check the finite computations their derivations use. The table below counts them by Part, and the index of the repository lists each with its grade, its premises and its Lean declarations.
\end{chtable}
\chstanding{instrument}
\chdeps{Definitions 1.1 to 1.6}{no theorem beyond it at this release}
\begin{chbridge}
The tally is an instrument of the compendium, and it draws no correspondence with the literature.
\end{chbridge}
\begin{chproof}[instrument]
The tally by Part. Each row counts the boxed statements of a Part by grade, with the theorems tied to named Lean declarations beside the theorems.
{\small
\begin{longtable}{lrrrrrr}
\hline
Part & theorems & in Lean & conditional & open results & tables & definitions \\
\hline
\endhead
\atlasByPartRows
\hline
\end{longtable}
}
\emph{Check.} The counts are generated from the boxed statements of the source by the index tool of the repository, whose checks regenerate them and compare them with the committed ones.
\end{chproof}
\label{chend:XVI-1-the-results-table-re-tallied-at-the-cut}
\par\medskip\noindent{\itshape The results tallied, the next chapter counts the open ones.}\label{narr:close:14-1}\par
\section{The open results census in physics terms}\label{mod:XIV-2-walls-census-in-physics-terms}
\chdesig{16-2}{14-2}{XIV-2-walls-census-in-physics-terms}{XIV-2-walls-census-in-physics-terms}
\label{narr:open:14-2}Table~238.1 prints the census of the open results by form, status and Part, the table to consult for every box of this compendium that is an open result.
\begin{chtable}[The census of the open results by form, status and Part]\label{box:XIV-2-walls-census-in-physics-terms:1}
The census counts 257 open results (the prior index of 178 plus 79 beyond it), equal at both of its engines; by form, 13 fit no form, 81 are a uniqueness claim where the construction admits several readings, 60 a number read at a single link, network or body, 26 a configuration beyond the reach of the closure at the dock rule, 35 an exploratory first parting of a closure and 43 a statement at the library rather than at the stable configurations (9 at a second form beside); by status 130 are superseded by a theorem established since, 49 re-posed under another name and 79 open (10 at a second status beside); and by Part they stand at I 7, II 35, III 14, IV 77, V 18, VI 16, VII 47, VIII 12, IX 7, X 4, XI 10, XII 8 and XIII 3 of the thirteen Parts of the earlier arrangement.
\end{chtable}
\chstanding{instrument}
\chdeps{Definitions 1.1 to 1.6 and Open results 245.1 and 245.2, Theorem 245.3, Open results 245.4 to 245.8 and Theorem 245.9}{Table 239.1 and Table 240.1}
In physics terms the six forms are the ways a derivation stays open. The first fits no form. The second is a uniqueness claim where several readings are admitted. The third is a number read at the bare link, a library network or one body, where the established numbers come from saturated stable configurations or joined reads. The fourth is a configuration beyond the reach of the closure, no qutrit end opening at the dock reading. The fifth is an exploratory first parting answered at a subsequent saturation. The sixth is a statement at the library rather than at the stable configurations. The three statuses are standing, re-posed under another name, and dissolved by a theorem established since, the superseding result named at every one. The census and one further result placed at this chapter, cited by its number, are read at their grades.
\begin{chbridge}
The literature calls the one link a single maximally entangled pair of two qudits (exact), the closure the forward closure of a set under operations (exact), the sink a stable configuration (nearest), the dock a dangling edge at a terminal qubit (nearest), and the objects named the open questions the Standard Model leaves, the hierarchy of its scales, the smallness of the strong CP angle, the three generations and the cosmological constant \cite{Weinberg1996,Navas2024}, by name only. The one link stands at Section~\ref{mod:I-2-link-and-its-six-genres}.
\end{chbridge}
\begin{chproof}[instrument]
The premises are the Axiom and the six definitions of Section~\ref{mod:FRAMEWORK} (definitions) and every open result of the construction at its own chapter (open results, read at their faces). The chapter uses nothing of them beyond their statements, and the chain from the axiom to this chapter is three chapters deep.
The table is the census in full at its three readings. The rule that produced it reads every open result at its own statement, types it at one form by the six rules of the census (a second form beside where two apply), marks it superseded or re-posed by the superseding or re-posing result or leaves it open, and files it at its Part. The first table counts by form with the second form beside, the second by status, the third by Part.
\begin{center}
\begin{adjustbox}{max width=\linewidth}\begin{tabular}{lrr}
\hline
form & count & with a second form beside \\
\hline
fitting no form & 13 & 16 \\
a uniqueness claim where several readings are admitted & 81 & 82 \\
a number read at a single link, network or body & 60 & 56 \\
a configuration beyond the reach of the closure at the dock rule & 26 & 26 \\
an exploratory first parting of a closure & 35 & 36 \\
a statement at the library rather than at the stable configurations & 43 & 42 \\
\hline
\end{tabular}\end{adjustbox}
\qquad
\begin{adjustbox}{max width=\linewidth}\begin{tabular}{lr}
\hline
status & count \\
\hline
superseded by a theorem established since & 130 \\
re-posed under another name & 49 \\
open & 79 \\
\hline
\end{tabular}\end{adjustbox}
\end{center}
\begin{center}
\begin{adjustbox}{max width=\linewidth}\begin{tabular}{lrlr}
\hline
Part & count & Part & count \\
\hline
I, the framework & 7 & VIII, decays and nuclear physics & 12 \\
II, quantum mechanics & 35 & IX, gravity and cosmology & 7 \\
III, spacetime and symmetry & 14 & X, thermodynamics and statistical physics & 4 \\
IV, gauge symmetry and the Standard Model & 77 & XI, condensed matter and atoms & 10 \\
V, quantum electrodynamics & 18 & XII, quantum information and computation & 8 \\
VI, chromodynamics and hadrons & 16 & XIII, the scorecard, the open results and what remains & 3 \\
VII, masses, mixing and generations & 47 & & \\
\hline
\end{tabular}\end{adjustbox}
\end{center}
The Parts are those of the earlier arrangement, thirteen in number. The counts sum to 258 at the token of the census, the 257 open results and one further line fitting no form.
At this release, beyond that census, six open results of this compendium close at the derivations that re-posed them, four as theorems and two as conditional theorems (the order register at the three directions, spin at the differing body, the probe read at the second route, the released body as the nested cross, beta decay at the literal closure reading and the decay path at the first no-path decay), twelve are re-posed at a new parting named, and one rule clause is struck beside the depth rule, each re-graded at its own chapter and re-tallied at the first chapter of this Part.
The check is that the two engines of the census agree at every count, that the index of 178 is re-met whole with the 79 beyond it complete, and that no figure is read against an observed value at this chapter.
The census was reproduced by an independent route, the two engines agreeing at every form, status and Part.
\end{chproof}

\par\medskip\noindent{\itshape The next chapter measures how much of the vocabulary of physics the construction has words for.}\label{narr:close:14-2}\par
\section{What remains: the ungraded page}\label{mod:XIV-7-what-remains-the-ungraded-page}
\chdesig{16-7}{14-7}{XIV-7-what-remains-the-ungraded-page}{XIV-7-what-remains-the-ungraded-page}
\label{narr:open:14-7}Table~239.1 prints the register of what the construction has not reached, and Open result~207.1, the strong reading as a theorem of the rate law, stands open.
\begin{chtable}[The register of what the construction has not reached]\label{box:XIV-7-what-remains-the-ungraded-page:1}
This chapter holds a register and states no result of the construction. The register is filled from the open namespace census of the seventh Lean release, decided in the kernel at its census theorem: 474 records, of them 74 open results at the census status open and 48 at re-posed (the 128 superseded and the 82 beyond the census not listed here), 133 conditional clauses of the theorems and 2 reading clauses of the twentieth atom of the definition ledger, every entry at its label, its status and its chapter as the census decides it; the entries are listed in the reference list of this chapter, each at its title and its number.
\end{chtable}
\chstanding{instrument}
\chdeps{Definitions 1.1 to 1.6 and Table 238.1}{Open result 207.1}
The 74 open results stand at their own statements. The 48 re-posed ones stand as the question first posed, a result beyond posing the same question under another name. The 133 conditional clauses are the premises at which the theorems of the construction are conditional, each declared as a proposition with no proof term. The 2 reading clauses are the two readings of the count of a read owed at its definition. What stands open is the register itself, every entry at its own grade, to be superseded at the next composition by the open-problems register and the derivation graph of the compendium, with the open results at the open results chapter; two further results placed at this chapter, cited by their numbers, are read at their grades.
\begin{chbridge}
The literature calls what a theory has not derived its open problems, the parameters it takes as inputs and the questions it leaves, the cosmological constant among them \cite{Navas2024,Weinberg1989}, and an ungraded page the list by domain of every entry a program has not reached with the gate each waits on, by name only.
\end{chbridge}
\begin{chproof}[instrument]
The premises are the Axiom and the six definitions of Section~\ref{mod:FRAMEWORK} (definitions) and every result the register lists at its own chapter. The chapter uses nothing of them beyond their statements, and the chain from the axiom to this chapter is three chapters deep.
The table is the register at its counts. The rule that produced it reads the open namespace of the seventh Lean release at its census theorem, decided in the kernel, and files every record at its status. The columns are the kind of entry and its count, the entries themselves listed at the reference list of this chapter, each at its title and its number.
\begin{center}
\begin{adjustbox}{max width=\linewidth}\begin{tabular}{lr}
\hline
the kind of entry at the open namespace & count \\
\hline
open results at the status open & 74 \\
open results re-posed under another name & 48 \\
conditional clauses of the theorems, each a proposition with no proof term & 133 \\
reading clauses of the twentieth atom of the definition ledger & 2 \\
superseded results, not listed here & 128 \\
results beyond the census, not listed here & 82 \\
\hline
the records of the census & 474 \\
\hline
\end{tabular}\end{adjustbox}
\end{center}
The check is that the six counts of the table sum to 467 of the 474 records of the census, that every entry stands at its own grade, and that no figure is read against an observed value at this chapter.
The register was reproduced by an independent route, the census theorem of the seventh Lean release decided in the kernel. At this release six of the open results of this compendium stand closed at theorems or conditional theorems (the order register at the three directions, spin at the differing body, the probe read at the second route, the released body as the nested cross, beta decay at the literal closure reading and the decay path at the first no-path decay) and twelve stand re-posed at a new parting, as the first chapter of this Part tallies them and their own chapters state them; the register above is carried as it stood.
\end{chproof}

\par\medskip\noindent{\itshape The correspondence sentences follow, one for every name the compendium uses.}\label{narr:close:14-7}\par
\section{Conclusion}\label{mod:XIV-C-conclusion}
\chdesig{16-C}{14-C}{XIV-C-conclusion}{XIV-C-conclusion}
\label{narr:open:14-C}Table~240.1 holds the conclusion, resting on the results table and the census of the open results.
\begin{chtable}[The conclusion, held]\label{box:XIV-C-conclusion:1}
This chapter holds the place of the conclusion and states no result of the construction.
\end{chtable}
\chstanding{instrument}
\chdeps{Definitions 1.1 to 1.6 and Table 238.1}{no theorem beyond it at this release}
What the chapter waits on is the conclusion itself, written from the completed Parts, with the results table and the census of the open results of this Part as its material.
\begin{chbridge}
The literature judges a theory of fundamental physics by what it derives from how little, the Standard Model with general relativity accounting for the observations at the price of its parameters and its unanswered questions \cite{Navas2024,Weinberg1995}, and a conclusion states what a program claims, at what status, and what it does not claim, by name only. This chapter draws no correspondence with the literature.
\end{chbridge}
\begin{chproof}[instrument]
The premises are the Axiom and the six definitions of Section~\ref{mod:FRAMEWORK} (definitions). The chapter uses nothing beyond their statements, and the chain from the axiom to this chapter is five chapters deep through the chapters of this Part.
The table is the placement of this chapter in the compendium. The rule that produced it places every established result at one chapter by its reading and marks a chapter at which no result is placed as held. The one column lists the material of the conclusion and where it stands.
\begin{center}
\begin{adjustbox}{max width=\linewidth}\begin{tabular}{ll}
\hline
the material of the conclusion & where it stands \\
\hline
the scorecard & the first chapter of this Part \\
the open results & the second chapter of this Part \\
\hline
\end{tabular}\end{adjustbox}
\end{center}
The check is that no figure is read against an observed value at this chapter.
The placement was reproduced by an independent route, the two engines agreeing at every chapter of the frame. At this release the material of the conclusion moves at the open results re-graded at their chapters, six closed (the order register at the three directions, spin at the differing body, the probe read at the second route, the released body as the nested cross, beta decay at the literal closure reading and the decay path at the first no-path decay) and twelve re-posed, tallied at the first chapter of this Part; the conclusion itself is held as above.
\end{chproof}
\par\medskip\noindent{\itshape Every theorem of this compendium descends from the one sentence at its start, and the reader who wants to see the chain for any one of them will find it in the dependency line beneath its box and in the graph at the front.}\label{narr:close:14-C}\par
